% =============================================================================
\section{Energía de un sistema de conductores}\label{sec:conductores}
% =============================================================================

Un conductor en equilibrio es una distribución superficial de carga con una peculiaridad: su superficie es equipotencial~\cite{T3}. Eso
simplifica enormemente la energía.

\begin{proposicion}[Energía de un sistema de conductores]\label{prop:conductores}
La energía electrostática de un sistema de $N$ conductores con cargas $Q_i$ y potenciales $\phi_i$ es
\begin{equation}\label{eq:conductores}
  \boxed{\;U_e=\frac12\sum_{i=1}^NQ_i\,\phi_i\;}
  =\frac12\sum_{i,j=1}^Np_{ij}\,Q_iQ_j=\frac12\sum_{i,j=1}^Nc_{ij}\,\phi_i\phi_j .
\end{equation}
\end{proposicion}

\begin{proof}
La carga reside en las superficies $S_i$ de los conductores, y sobre cada una el potencial es constante, $\phi=\phi_i$. Por la
\cref{prop:continua},
\[
  U_e=\frac12\sum_{i=1}^N\oiint_{S_i}\sigma_i\,\phi\,\dd S=\frac12\sum_{i=1}^N\phi_i\oiint_{S_i}\sigma_i\,\dd S=\frac12\sum_{i=1}^N\phi_i\,Q_i .
\]
Las otras dos formas se obtienen sustituyendo $\phi_i=\sum_jp_{ij}Q_j$ o $Q_i=\sum_jc_{ij}\phi_j$~\cite{T3}.
\end{proof}

\begin{observacion}[Lectura de la fórmula]\label{obs:conductores}
\begin{enumerate}
  \item Formalmente se parece a la energía de cargas puntuales \eqref{eq:puntuales}, pero el significado es otro: aquí $\phi_i$ es el
        potencial del conductor $i$ creado por \emph{todas} las cargas, incluida la suya, $Q_i$. Es energía \emph{total}, no solo de
        interacción, y por el \cref{cor:positiva} es positiva o nula.
  \item En consecuencia, las formas cuadráticas $\sum p_{ij}Q_iQ_j$ y $\sum c_{ij}\phi_i\phi_j$ son definidas positivas: se anulan solo si todas las
        cargas (o todos los potenciales) son nulos. Es otra forma de ver las propiedades de signo de $p_{ij}$ y $c_{ij}$ del Tema~3.
  \item Un conductor a tierra ($\phi_i=0$) o descargado ($Q_i=0$) no contribuye explícitamente a la suma, aunque sí a través de los
        potenciales o las cargas de los demás.
\end{enumerate}
\end{observacion}

\begin{corolario}[Energía de un condensador]\label{cor:condensador}
Un condensador de capacidad $C$, con cargas $\pm Q$ en sus armaduras y diferencia de potencial $V=\phi_1-\phi_2$, almacena la energía
\begin{equation}\label{eq:condensador}
  \boxed{\;U_e=\frac12\,QV=\frac12\,CV^2=\frac{Q^2}{2C}\;}.
\end{equation}
\end{corolario}

\begin{proof}
Con $Q_1=Q$ y $Q_2=-Q$, la \cref{prop:conductores} da $U_e=\frac12\big[Q\phi_1+(-Q)\phi_2\big]=\frac12Q(\phi_1-\phi_2)=\frac12QV$, y las otras dos formas se
siguen de $Q=CV$.
\end{proof}

\begin{ejemplo}[Esfera conductora aislada]\label{ej:esfera-aislada}
Para una esfera de radio $a$ y carga $Q$, $C=4\pi\eo a$ (Tema~3) y $U_e=Q^2/2C=Q^2/(8\pi\eo a)$, como en el \cref{ej:corteza}. Con $p_{11}=1/C$, es
también $\frac12p_{11}Q^2$.
\end{ejemplo}

\begin{ejemplo}[Condensador plano]\label{ej:plano}
Con $C=\eo S/d$, la energía es $U_e=\frac12CV^2=\frac12\eo SV^2/d$. Con el campo, $E=V/d$ uniforme en el volumen $Sd$ entre las placas y nulo
fuera (despreciando los bordes): $U_e=\frac{\eo}{2}\big(V/d\big)^2Sd$, idéntico. La energía está repartida uniformemente entre las placas.
\end{ejemplo}
